\subsection{用展开式定义实数}

反之，设给定了一个整数 \(q_0\) 与整数的序列 \((v_n)\) ( \(n \geqslant 1\) )（满足 \(0 \leqslant v_n \leqslant a_n - 1\) ）；我们问是否存在一个数 \(x\) ，其展开 (4) 满足 \(p_0 = q_0\) 且 \(u_n = v_n\) 对一切 \(n\) 。如果这样的数存在，它就是唯一的，因为它等于 \(q_0 + \sum_{n=1}^{\infty} \frac{v_n}{d_n}\) 。

对每个整数 \(m > 0\) 我们有（由比较原理）

\[
\sum_ {n = m + 1} ^ {\infty} \frac {v _ {n}}{d _ {n}} \leqslant \sum_ {n = m + 1} ^ {\infty} \frac {a _ {n} - 1}{d _ {n}} = \sum_ {n = m + 1} ^ {\infty} \left(\frac {1}{d _ {n - 1}} - \frac {1}{d _ {n}}\right) = \frac {1}{d _ {m}}
\]

而且最左端与最右端的项仅当 \(v_{n} = a_{n} - 1\) （对每个 \(n > m\) ）时才相等（§ 7，第 1 节，定理 2）。因此通项为 \(\frac{v_n}{d_n}\) 的级数收敛；此外，若 \(x = q_0 + \sum_{n=1}^{\infty} \frac{v_n}{d_n}\) ，我们有

\[
s _ {m} = q _ {0} + \sum_ {n = 1} ^ {m} \frac {v _ {n}}{d _ {n}} \leqslant x \leqslant s _ {m} + \frac {1}{d _ {m}}
\]

且 \(x = s_m + \frac{1}{d_m}\) 仅当 \(v_n = a_n - 1\) （对每个 \(n > m\) ）时成立。由于 \(s_m\) 是分母为 \(d_m\) 的分数，不足近似 \(r_m\) （它对 \(x\) 精确到 \(1/d_m\) ）等于 \(s_m\) 或 \(s_m + \frac{1}{d_m}\) ；且后一情形仅当 \(v_{n} = a_{n} - 1\) （对一切 \(n > m\) ）时才能发生。于是：

\begin{enumerate}
\def\labelenumi{(\roman{enumi})}
\item
  存在无穷多个 \(n\) 的值（满足 \(v_{n} < a_{n} - 1\) ）：这时级数 \(q_{0} + \sum_{n=1}^{\infty} \frac{v_{n}}{d_{n}}\) 与其和 \(x\) 的展开完全相同。
\item
  存在一个整数 \(m \geqslant 0\) 使得 \(v_{n} = a_{n} - 1\) （只要 \(n > m\) ），且 \(v_{m} < a_{m} - 1\) （若 \(m > 0\) ）；这时 \(x\) （级数 \(q_{0} + \sum_{n=1}^{\infty} \frac{v_{n}}{d_{n}}\) 的和）等于有理数
\end{enumerate}

\[
q _ {0} + \sum_ {n = 1} ^ {m} \frac {v _ {n}}{d _ {n}} + \frac {1}{d _ {m}}\tag{5}
\]

它是 \(k/d_{m}\) 型的（k 为整数）；x 的展开与级数 (5) 完全相同，其中所有指标 \textgreater{} m 的项都是零；这样的展开称为是有限的，或说是终止的。级数

\[
q _ {0} + \sum_ {n = 1} ^ {\infty} \frac {v _ {n}}{d _ {n}} = q _ {0} + \sum_ {n = 1} ^ {m} \frac {v _ {n}}{d _ {n}} + \sum_ {n = m + 1} ^ {\infty} \frac {a _ {n} - 1}{d _ {n}}\tag{6}
\]

称为数 x 的非正常展开。

反之，设 x 是一个有理数，它对 n 的某个值可以写成分母为 \(d_{n}\) 的分数。~设 m 是使 x 为 \(k/d_{m}\) 型（k 为整数）的最小整数；我们有 \(r_{n}<x\) （对 n\textless m），且 \(r_{m}=x\) ，因此 x 的展开具有形式 (5)，且 x 有一个由 (6) 给出的非正常展开；此外这个非正常展开是唯一的。

一个有理数，写成既约形式 \(p / q\) ，等于分母为 \(d_n\) 的分数当且仅当 \(q\) 整除 \(d_n\) （数 \(m\) 因而是这样的最小整数 \(n\) ：\(q\) 整除 \(d_n\) ）。可能发生这样的情形：每个有理数都具有这一性质（对适当选取的 \(n\) ）；当且仅当每个整数 \(> 0\) 整除某个 \(d_n\) （例如 \(d_n = n!\) ）时才是这种情形。如果诸 \(d_n\) 具有这一性质，那么一个数是有理数当且仅当它关于序列 \((d_n)\) 的展开终止。

概括起来：对每个序列 \(s\) ，其首项 \(q_{0}\) 是任意整数，且其项 \(v_{n}(n \geqslant 1)\) 满足 \(0 \leqslant v_{n} \leqslant a_{n} - 1\) ，对应着一个等于 \(q_{0} + \sum_{n=1}^{\infty} \frac{v_{n}}{d_{n}}\) 的实数；若 \(I_{n}\) 表示区间 \([0, a_{n} - 1]\) （\(N\) 中的），我们就这样定义了一个映射 \(\varphi\) （从 \(X = Z \times \prod_{n=1}^{\infty} I_{n}\) 到直线 \(R\) 上）；此外，方程 \(\varphi(s) = x\) ——其中 \(x \in R\) 已给定——当 \(x\) 不是分母为 \(d_{n}\) （对某个 \(n\) ）的分数时有一个解，否则有两个解。

